% SGA 3, Expose VI_B, Loop 1 English reconstruction.
% Scope: Definition 7.2.2, Recall 7.3.0, and Corollary 7.3
% with its complete proof, including source editor's notes 73--74.
% Authority: Polo--Gille Exp6B-13oct24.pdf, component-local pages
% 52--54 (combined-reader pages 393--395; re-edition headers
% 382--384; original-margin printed pages 365--367).
% Proposition 7.4 begins on local page 54 and is excluded.
% OCR and Reinhold English Markdown are comparison witnesses only.
% This unit contains no diagrams.

\medskip
\phantomsection\label{def:sga3-VIB-7.2.2}%
\noindent\textbf{定义 7.2.2.---}
设 \(G\) 为 \(k\)-群.
\begin{enumerate}
\item[\textup{a)}]\phantomsection\label{def:sga3-VIB-7.2.2-a}
  设 \(A\) 和 \(B\) 为 \(G\) 的两个几何约化的
  \(k\)-子群概形. \emph{\(A\) 和 \(B\) 在 \(G\) 中的换位子群
  概形}, 记为
  \[
    (A,B)_G
    \qquad\text{或简记为}\qquad
    (A,B),
  \]
  是 \(G\) 中由态射
  \phantomsection\label{form:sga3-VIB-7.2.2-commutator}%
  \[
    \nu:A\times_kB\longrightarrow G
  \]
  生成的闭子群概形; 此态射对每个 \(k\)-概形 \(S\) 定义为
  \[
    (a,b)\longmapsto aba^{-1}b^{-1},
    \qquad
    a\in A(S),\quad b\in B(S).
  \]

\item[\textup{b)}]\phantomsection\label{def:sga3-VIB-7.2.2-b}
  假设 \(G\) 在 \(k\) 上几何约化.
  \emph{\(G\) 的导群}, 记为 \(\mathfrak{D}(G)\),%
  \footnote{\label{note:sga3-VIB-73}编者注: 将原文的记号
  \(\operatorname{Der}(G)\) 改掉,
  以免与导子空间混淆.}
  是群
  \[
    \mathfrak{D}(G)=(G,G).
  \]
\end{enumerate}

\noindent\emph{注意.}
群 \(G\) 交换当且仅当 \(\mathfrak{D}(G)\) 是
单位 \(k\)-群.

\medskip
\phantomsection\label{recall:sga3-VIB-7.3.0}%
\noindent\textbf{回顾 \SGARefLink{recall:sga3-VIB-7.3.0}{7.3.0}.---}%
\footnote{\label{note:sga3-VIB-74}编者注: 这些回顾
为新增内容, 并相应地修改和扩充了
\SGARefLink{cor:sga3-VIB-7.3}{7.3} 的陈述.}
设 \(\varphi:Y\to Z\) 为 \(S\)-概形的态射. 回顾,
像预层 \(\varphi(Y)\) 是如下 \(S\)-函子:
它将 \(S\) 上的每个 \(S'\) 映到子集
\[
  \varphi\bigl(Y(S')\bigr)\subset Z(S').
\]
假设 \(T\) 为 \(Z\) 的子概形. 若预层的包含态射
\[
  \varphi(Y)\hookrightarrow Z
\]
经由 \(T\) 分解, 即若对每个 \(S\)-态射 \(h:S'\to Y\),
都有
\[
  \varphi\circ h\in T(S')
\]
则在集合意义下,
\(\varphi(Y)\subset T\): 只需取 \(h=\operatorname{id}_Y\).
反之, 若 \(Y\) 约化, 且在集合意义下
\(\varphi(Y)\subset T\), 则 \(\varphi\)
经由 \(T\) 分解.

还要回顾, 由 \hyperref[prop:sga3-VIB-6.2]{6.2}, 若
\(H\) 为 \(G\) 的闭 \(k\)-子群概形, 则
\[
  \underline{\operatorname{Centr}}_G H
  \qquad\text{和}\qquad
  \underline{\operatorname{Norm}}_G H
\]
可由 \(G\) 的闭 \(k\)-子群概形表出.

\medskip
\phantomsection\label{cor:sga3-VIB-7.3}%
\noindent\textbf{推论 7.3.---}
\emph{设 \(G\) 为 \(k\)-群, \(X\) 为几何
约化的 \(k\)-概形, \(f:X\to G\) 为
\(k\)-态射.}
\begin{enumerate}
\item[\textup{\SGAInlineRefTarget{cor:sga3-VIB-7.3:item:i}(i)}]\phantomsection\label{concl:sga3-VIB-7.3-i}
  \emph{设 \(S\) 为 \(k\)-概形, \(u\) 为
  \(S\)-群 \(G_S\) 的自态射.}
  \begin{enumerate}
  \item[\textup{(a)}]\phantomsection\label{concl:sga3-VIB-7.3-i-a}
    \emph{若在集合意义下}
    \[
      u\bigl(f_S(X_S)\bigr)\subset\Gamma_G(f)_S
    \]
    \emph{成立, 则态射}
    \[
      u:\Gamma_G(f)_S\longrightarrow G_S
    \]
    \emph{经由 \(\Gamma_G(f)_S\) 分解.}

  \item[\textup{\SGAInlineRefTarget{cor:sga3-VIB-7.3:item:b}(b)}]\phantomsection\label{concl:sga3-VIB-7.3-i-b}
    \emph{若 \(u\) 为 \(S\)-群 \(G_S\) 的自同构,
    且在集合意义下}
    \[
      u\bigl(f_S(X_S)\bigr)=f_S(X_S)
    \]
    \emph{成立, 则 \(u\) 诱导 \(\Gamma_G(f)_S\)
    的自同构. 特别地, 若 \(g\in G(S)\)
    满足在集合意义下}
    \[
      \operatorname{int}(g)\bigl(f_S(X_S)\bigr)=f_S(X_S)
    \]
    \emph{成立, 则}
    \[
      g\in
      \underline{\operatorname{Norm}}_{G_S}
        \bigl(\Gamma_G(f)_S\bigr)(S).
    \]

  \item[\textup{\SGAInlineRefTarget{cor:sga3-VIB-7.3:item:c}(c)}]\phantomsection\label{concl:sga3-VIB-7.3-i-c}
    \emph{若 \(u\circ f_S=f_S\), 则 \(u\) 在
    \(G_S\) 的子群 \(\Gamma_G(f)_S\) 上的限制为恒等态射.
    特别地, 若 \(g\in G(S)\) 满足}
    \[
      \operatorname{int}(g)\circ f_S=f_S,
    \]
    \emph{则}
    \[
      g\in
      \underline{\operatorname{Centr}}_{G_S}
        \bigl(\Gamma_G(f)_S\bigr)(S).
    \]
  \end{enumerate}

\item[\textup{\SGAInlineRefTarget{cor:sga3-VIB-7.3:item:ii}(ii)}]\phantomsection\label{concl:sga3-VIB-7.3-ii}
  \emph{设 \(H\) 为 \(G\) 的子群概形. 则 \(H\)
  中心化 (相应地, 正规化) \(\Gamma_G(f)\) 当且
  仅当对每个 \(k\)-概形 \(S\) 及每个 \(h\in H(S)\),}
  \[
    \operatorname{int}(h)\circ f_S=f_S
  \]
  \emph{(相应地, 在集合意义下}
  \[
    \operatorname{int}(h)\bigl(f_S(X_S)\bigr)
      \subset\Gamma_G(f)_S
  \]
  \emph{成立). 等价地, 对每个
  \(S'\to S\) 及每个 \(x\in X(S')\), \(G(S')\) 中的元素 \(h_{S'}\)
  和 \(f(x)\) 交换 (相应地,}
  \[
    h_{S'}f(x)h_{S'}^{-1}\in\Gamma_G(f)(S').
  \]
  \emph{).}

\item[\textup{\SGAInlineRefTarget{cor:sga3-VIB-7.3:item:iii}(iii)}]\phantomsection\label{concl:sga3-VIB-7.3-iii}
  \emph{特别地, 设 \(Y\) 为另一个几何
  约化的 \(k\)-概形, \(\varphi:Y\to G\) 为
  \(k\)-态射. 假设对每个 \(k\)-概形 \(S'\),
  每个 \(x\in X(S')\) 及每个 \(y\in Y(S')\), \(G(S')\) 中的元素
  \(\varphi(y)\) 和 \(f(x)\) 交换
  (相应地,}
  \[
    \operatorname{int}\bigl(\varphi(y)\bigr)\bigl(f(x)\bigr)
      \in\Gamma_G(f)(S').
  \]
  \emph{).}
  \emph{则 \(\Gamma_G(\varphi)\) 为}
  \[
    \underline{\operatorname{Centr}}_G\bigl(\Gamma_G(f)\bigr)
    \qquad\text{\emph{(相应地, 为)}}\qquad
    \underline{\operatorname{Norm}}_G\bigl(\Gamma_G(f)\bigr).
  \]
  \emph{的 \(k\)-子群.}

\item[\textup{\SGAInlineRefTarget{cor:sga3-VIB-7.3:item:iv}(iv)}]\phantomsection\label{concl:sga3-VIB-7.3-iv}
  \emph{若 \(g\) 为 \(G\) 的 \(k\)-点, 则 \(k\)-群
  \(\Gamma_G(\{g\})\) 交换.}

\item[\textup{(v)}]\phantomsection\label{concl:sga3-VIB-7.3-v}
  \emph{设 \(A\) 和 \(B\) 为 \(G\) 的两个子群概形,
  它们在 \(k\) 上几何约化. 若 \(A\) 和 \(B\)
  正规 (相应地, 为特征的), 则
  \((A,B)\) 也有同一性质.}
\end{enumerate}

\emph{证明.}
先证明 \textup{\SGARefLink{cor:sga3-VIB-7.3:item:i}{(i)}\SGARefLink{concl:sga3-VIB-7.3-i}{(a)}}. 令
\[
  \Gamma_S=\Gamma_G(f)_S.
\]
则 \(u^{-1}(\Gamma_S)\) 为 \(G_S\) 的闭 \(S\)-子群概形.
因此由
\hyperref[concl:sga3-VIB-7.1-iii]{7.1\textup{(iii)}},
只需证明 \(f_S\) 经由
\(u^{-1}(\Gamma_S)\) 分解. 由于
\[
  u\bigl(f_S(X_S)\bigr)\subset\Gamma_S
\]
且 \(X_S\) 约化, \(u\circ f_S\) 经由
\(\Gamma_S\) 分解; 从而 \(f_S\) 经由
\(u^{-1}(\Gamma_S)\) 分解. 这就证明了 \textup{\SGARefLink{concl:sga3-VIB-7.3-i-a}{(a)}}.

将 \textup{\SGARefLink{concl:sga3-VIB-7.3-i-a}{(a)}} 分别应用于 \(u\) 和 \(u^{-1}\),
即得 \textup{\SGARefLink{cor:sga3-VIB-7.3:item:b}{(b)}} 的第一个断言. 事实上, 只需假设
\[
  u\bigl(f_S(X_S)\bigr)\subset\Gamma_S
  \qquad\text{且}\qquad
  u^{-1}\bigl(f_S(X_S)\bigr)\subset\Gamma_S.
\]
第二个断言为 \(u=\operatorname{int}(g)\) 的特殊情形.

现在证明 \textup{\SGARefLink{concl:sga3-VIB-7.3-i-c}{(c)}}. 考虑 \(S\)-群态射
\[
  G_S\longrightarrow G_S\times_SG_S,
  \qquad
  g\longmapsto\bigl(u(g),g\bigr),
\]
并设 \(H\) 为对角线的逆像. 它是 \(G_S\) 的子群
概形. 由于 \(G\) 在 \(k\) 上分离, 见
\hyperref[sga3:VIA:0.3]{VI\(_A\), 0.3}, \(G_S\) 在
\(S\) 上分离, 因而 \(H\) 在 \(G_S\) 中闭. 由假设,
\(f_S\) 经由 \(H\) 分解; 因此由
\hyperref[concl:sga3-VIB-7.1-iii]{7.1\textup{(iii)}}, \(H\) 包含
\(\Gamma_G(f)_S\). 这就证明了
\textup{\SGARefLink{cor:sga3-VIB-7.3:item:c}{(c)}} 的第一个断言, 第二个断言为
\(u=\operatorname{int}(g)\) 的特殊情形.

接着证明 \textup{\SGARefLink{cor:sga3-VIB-7.3:item:ii}{(ii)}}. 令
\[
  \Gamma=\Gamma_G(f)
\]
并用 \(i:\Gamma\hookrightarrow G\) 表示包含态射.
群 \(H\) 包含于
\[
  N=\underline{\operatorname{Norm}}_G(\Gamma)
\]
当且仅当对每个 \(k\)-概形 \(S\) 及每个
\(h\in H(S)\),
\[
  \operatorname{int}(h)(\Gamma_S)\subset\Gamma_S.
\]
将此条件应用于 \(h\) 和 \(h^{-1}\), 可知此时该
包含关系为等式. 由 \textup{\SGARefLink{cor:sga3-VIB-7.3:item:i}{(i)}\SGARefLink{concl:sga3-VIB-7.3-i}{(a)}}, 此条件
成立当且仅当
\[
  \operatorname{int}(h)\bigl(f_S(X_S)\bigr)\subset\Gamma_S.
\]
同样, \(H\) 包含于
\[
  C=\underline{\operatorname{Centr}}_G(\Gamma)
\]
当且仅当对每个 \(k\)-概形 \(S\) 及每个
\(h\in H(S)\),
\[
  \operatorname{int}(h)\circ i_S=i_S.
\]
由 \textup{\SGARefLink{cor:sga3-VIB-7.3:item:i}{(i)}\SGARefLink{cor:sga3-VIB-7.3:item:c}{(c)}}, 这成立当且仅当
\[
  \operatorname{int}(h)\circ f_S=f_S.
\]
这就证明了 \textup{\SGARefLink{cor:sga3-VIB-7.3:item:ii}{(ii)}}.

现在证明 \textup{\SGARefLink{cor:sga3-VIB-7.3:item:iii}{(iii)}}. 根据 \textup{\SGARefLink{cor:sga3-VIB-7.3:item:ii}{(ii)}},
假设蕴涵 \(\varphi\) 经由 \(C\)
(相应地, \(N\)) 分解. 后者由 \hyperref[prop:sga3-VIB-6.2]{6.2}
为 \(G\) 的闭子群, 因而由
\hyperref[concl:sga3-VIB-7.1-i]{7.1\textup{(i)}}, 包含
\(\Gamma_G(\varphi)\).%
\begingroup
\renewcommand{\thefootnote}{\fnsymbol{footnote}}%
\footnote[1]{\label{note:sga3-VIB-7.3-Gamma-phi}译者注:
法文底本此处印作 \(\Gamma_G(f)\). 由于这里证明经由
\(C\) (相应地, \(N\)) 分解的态射为
\(\varphi\), 而非 \(f\), 所以 \SGARefLink{prop:sga3-VIB-7.1}{7.1}\textup{\SGARefLink{concl:sga3-VIB-7.1-i}{(i)}} 给出
\(\Gamma_G(\varphi)\), 与陈述一致.}
\endgroup

在 \textup{\SGARefLink{cor:sga3-VIB-7.3:item:iii}{(iii)}} 中令 \(f\) 和 \(\varphi\)
均为 \(g\) 定义的闭浸入
\[
  \operatorname{Spec}k\hookrightarrow G
\]
即得到断言 \textup{\SGARefLink{cor:sga3-VIB-7.3:item:iv}{(iv)}}.
事实上, 对每个 \(k\)-概形 \(S\), 此时集合
\(X(S)\) 和 \(Y(S)\) 都仅有一个元素,
分别由 \(f\) 和 \(\varphi\) 映到 \(g\).

最后证明 \textup{\SGARefLink{concl:sga3-VIB-7.3-iv}{(v)}}. 设
\[
  \nu:G\times_kG\longrightarrow G,
  \qquad
  (g,h)\longmapsto ghg^{-1}h^{-1},
\]
并设 \(\nu'\) 为其在 \(A\times_kB\) 上的限制. 由
\hyperref[def:sga3-VIB-7.2.2]{7.2.2},
\[
  (A,B)=\Gamma_G(\nu').
\]
设 \(S\) 为 \(k\)-概形, \(u\) 为 \(G_S\) 的一个内
自同构, 相应地, 任意一个 \(S\)-群
自同构. 则
\[
  u\circ\nu_S=\nu_S\circ(u\times_Su).
\]
由假设, \(u\) 诱导 \(A_S\) 的自同构 \(u_1\)
及 \(B_S\) 的自同构 \(u_2\). 因此,
\phantomsection\label{form:sga3-VIB-7.3-v-commutator-image}%
\[
  u\bigl(\nu'_S(A_S\times_SB_S)\bigr)
    =\nu'_S\bigl(u_1(A_S)\times_Su_2(B_S)\bigr)
    =\nu'_S(A_S\times_SB_S).
\]
所以由 \textup{\SGARefLink{concl:sga3-VIB-7.3-i}{(i)}\SGARefLink{concl:sga3-VIB-7.3-i-b}{(b)}}, \(u\) 诱导
\((A,B)_S\) 的自同构. 这就证明了 \textup{\SGARefLink{concl:sga3-VIB-7.3-v}{(v)}}.

% 翻译统计: 源行数 324; 输出行数 326 (含本统计行); 行内/行间数学片段 212; 翻译占位 0. 底本疑点另报主线程.

% 独立复核检查点：已逐段对照英文源文，覆盖正文、假设、证明、脚注与源注。
% 行内/陈列数学及全部引用键机械对照通过；仅统一群的invariant为“正规”，保留源文疑点另报。
